% ==========================================================================
%  SPHERExLens -- presentation specification: every slide, every figure and
%  the complete spoken script. 30 slides; full route 28:10, 15-minute route 18:40.
%  Layout follows the reference "slides and script" document:
%  title page + deck-map table, then one page per slide with
%  "On the slide." / "Figure / diagram." / the figure / "Spoken script."
%  Build:  pdflatex slides_and_script.tex   (twice)
% ==========================================================================
\documentclass[11pt,a4paper]{article}

\usepackage[T1]{fontenc}
\usepackage[utf8]{inputenc}
\usepackage{lmodern}
\usepackage{microtype}
\usepackage[a4paper,margin=2.0cm,top=2.0cm,bottom=2.0cm]{geometry}
\usepackage{amsmath,amssymb}
\usepackage{booktabs,array,longtable,multirow}
\usepackage{graphicx}
\usepackage{xcolor}
\usepackage{enumitem}
\usepackage{titlesec}
\usepackage{fancyhdr}
\usepackage{tikz}
\usetikzlibrary{arrows.meta,positioning,shapes.geometric,calc,fit,backgrounds,patterns,decorations.pathreplacing}
\usepackage{pgfplots}
\pgfplotsset{compat=1.18}
\usepackage{listings}
\usepackage[hidelinks]{hyperref}

% --------------------------------------------------------------- colours ---
\definecolor{ink}{HTML}{14213A}
\definecolor{slate}{HTML}{3F4A5A}
\definecolor{muted}{HTML}{7B8798}
\definecolor{accent}{HTML}{1F6FB2}
\definecolor{accent2}{HTML}{E8F1F8}
\definecolor{good}{HTML}{1E7B4F}
\definecolor{good2}{HTML}{E7F3EC}
\definecolor{bad}{HTML}{B3261E}
\definecolor{bad2}{HTML}{FAECEA}
\definecolor{warn}{HTML}{8A5A00}
\definecolor{warn2}{HTML}{FBF3E2}
\definecolor{rule}{HTML}{C9D0DA}
\definecolor{purple}{HTML}{7C3AED}
\definecolor{purple2}{HTML}{F3E8FF}

% ------------------------------------------------------------ page style ---
\pagestyle{plain}
\setlength{\parindent}{0pt}
\setlength{\parskip}{4pt}

% ---------------------------------------------------------------- macros ---
\newcommand{\arcmin}{\ensuremath{{}^{\prime}}}
\newcommand{\arcsec}{\ensuremath{{}^{\prime\prime}}}
\newcommand{\code}[1]{\texttt{\small\detokenize{#1}}}

% slide header:  Slide 1 --- Title       0:45
% The header is measured first and scaled down to the text width when a long
% title would otherwise run into the running time on the right.
\newsavebox{\slideheadbox}
\newcommand{\slidehead}[3]{%
  \sbox{\slideheadbox}{\large\bfseries Slide #1 --- #2\hspace{1em plus 1fill}#3}%
  \noindent\ifdim\wd\slideheadbox>\linewidth
    \resizebox{\linewidth}{!}{\usebox{\slideheadbox}}%
  \else
    \makebox[\linewidth][s]{\usebox{\slideheadbox}}%
  \fi\par
  \vspace{2pt}\hrule height 0.9pt\vspace{8pt}}

% figure wrapper: sized to fit cleanly without overflowing margins
\newcommand{\figbox}[1]{\begin{center}\resizebox{0.98\linewidth}{!}{#1}\end{center}}


\usepackage[active,tightpage]{preview}
\PreviewEnvironment{tikzpicture}
\begin{document}
\begin{tikzpicture}[font=\footnotesize,
  k/.style={draw=good,thick,fill=good2,align=center,inner sep=3pt,text width=3.3cm,rounded corners=2pt},
  c/.style={draw=warn,thick,fill=warn2,align=center,inner sep=3pt,text width=3.3cm,rounded corners=2pt},
  m/.style={draw=accent,thick,fill=accent2,align=center,inner sep=3pt,text width=3.3cm,rounded corners=2pt},
  ar/.style={-{Stealth[length=5pt]},draw=slate,thick}]

  % Left panel: Astrometric Epoch Propagation Geometry
  \draw[rule,thick,fill=white,rounded corners=2pt] (0.2,1.7) rectangle (7.6,6.4);
  \node[anchor=north west,font=\scriptsize\bfseries,text=ink] at (0.4,6.2) {Proper-Motion Epoch Propagation};

  \fill[accent] (1.7,4.6) circle (0.12);
  \node[font=\tiny\bfseries,text=accent,below] at (1.7,4.45) {Gaia J2016.0};

  \fill[good] (5.9,4.6) circle (0.12);
  \node[font=\tiny\bfseries,text=good,below] at (5.9,4.45) {SPHEREx 2026.1};

  % Unpropagated cone (miss)
  \draw[bad,dashed,thick] (1.7,4.6) circle (0.8);
  \node[font=\tiny,text=bad,above] at (1.7,5.4) {Naive cone: MISS};

  % Propagated cone (hit)
  \draw[good,thick] (5.9,4.6) circle (0.8);
  \node[font=\tiny,text=good,above] at (5.9,5.4) {Propagated: HIT};

  \draw[-{Stealth[length=6pt]},thick,good] (1.8,4.8) -- node[above,font=\tiny,text=good] {$10.4\arcsec/\mathrm{yr} \times 9.8\,\mathrm{yr} \approx 102\arcsec$} (5.8,4.8);

  \node[font=\tiny,text=muted,align=left,text width=6.6cm] at (4.0,2.65) {
    \textbf{The Barnard's Star Bug Resolved:}\\
    At $10.4\arcsec/\mathrm{yr}$, Barnard's Star is displaced by $\sim 94\arcsec$ between Gaia DR3 and SPHEREx DR1. Naive matching missed it, branding the star an unexplained CANDIDATE. In \code{catalog_tools.py}, \code{propagate_catalogs()} moves coordinates to exposure MJD, identifying it cleanly as \code{KNOWN_MOVER}.
  };

  % Right panel: 6-Class Triage Tree
  \draw[rule,thick,fill=white,rounded corners=2pt] (8.2,1.45) rectangle (16.0,6.4);
  \node[anchor=north west,font=\scriptsize\bfseries,text=ink] at (8.4,6.2) {Six-Category Multi-Tier Triage Flow};

  \node[m,text width=3.8cm] (in) at (12.1,4.9) {\textbf{Detection Coordinate}\\ Propagate Gaia/AllWISE to MJD};

  \node[k,text width=3.3cm] (km) at (10.1,3.6) {\textbf{\code{KNOWN_MOVER}}\\ \tiny Gaia match + motion\\ \textbf{\code{VARIABLE}} \tiny (VSX)};
  \node[m,text width=3.3cm] (art) at (14.1,3.6) {\textbf{\code{BRIGHT_STAR}}\\ \tiny $G \le 7$, $W1 \le 5.5$ halo\\ \textbf{\code{MASK_ARTEFACT}}};
  \node[c,text width=4.0cm] (can) at (12.1,2.4) {\textbf{\code{CANDIDATE} (Unexplained)}\\ \tiny Survives all multi-wavelength vetoes};

  \draw[ar] (in) -| (km);
  \draw[ar] (in) -| (art);
  \draw[ar] (in) -- (can);

  \node[draw=good,thick,fill=good2,rounded corners=2pt,align=center,inner sep=3.5pt,text width=15.4cm,font=\scriptsize] at (8.1,0.65)
    {\textbf{Triage Verdict Performance Across 18 Fields:} 99 of 109 detections (91\%) explained by catalogued astrometry, stellar variability, or optical artefacts, leaving exactly 10 candidate residuals for scientific follow-up.};
\end{tikzpicture}
\end{document}
